\chapter{术语表 Glossary}

本附录统一整理全书使用的核心术语，并显式解决各讲之间最容易混淆的概念边界。凡与本附录冲突的局部写法，均以本附录为准。

\section{推理、训练与工作流}

\begin{description}
\item[\textbf{推理时计算（inference-time computation）}] 部署阶段为了单个任务额外投入的计算预算，包括更长的 reasoning trace、更宽的并行采样、更深的修订循环、verifier 调用或 tree search。它描述的是\textbf{预算如何花}，不是某一种具体算法。

\item[\textbf{推理（reasoning）}] 生成、检查、修订中间步骤以支持答案或动作的过程。它可以是显式 chain-of-thought，也可以由 verifier、judge 或 environment feedback 驱动。

\item[\textbf{规划（planning）}] 面向未来动作的显式分解、前瞻或 world-model 式决策。相较于 reasoning，planning 更强调多步行动结构，而不是单次回答内部的逻辑展开。

\item[\textbf{搜索（search）}] 对多个候选 reasoning path、action trajectory 或 proof branch 进行枚举、扩展与筛选。搜索是显式多分支过程，而不等于任何较长的单轨推理。

\item[\textbf{记忆（memory）}] 超出当前上下文窗口的状态保留与检索机制，包括长期记忆、外部 memory store、RAG 式知识库或项目级上下文。它不是所有历史 token 的统称，而是\textbf{可再访问的持久状态}。

\item[\textbf{后训练（post-training）}] 预训练之后、面向任务行为塑形的阶段。它是总类，下面可以包含 instruction tuning、preference tuning、RL fine-tuning 等不同做法。

\item[\textbf{基于人类反馈的强化学习（RLHF）}] 一类利用人类偏好或判断信号来训练策略模型的后训练家族。书中将 RLHF 视为框架，而不是某个唯一 loss。

\item[\textbf{直接偏好优化（DPO）}] 直接利用 preference pairs 对策略模型做优化、并通过 reference policy 保持更新稳定的方法。它是偏好优化器，不是全部 reasoning recipe。

\item[\textbf{近端策略优化（PPO）}] 一种常见的强化学习算法，常被放入 RLHF pipeline 中。书中只在具体 policy optimization 语境下使用该术语，不把它当成 RLHF 的同义词。

\item[\textbf{组相对策略优化（GRPO）}] 以一组候选样本的相对优劣为基础做策略更新的优化方法。使用时保留英文缩写，并说明比较是如何分组与归一化的。

\item[\textbf{可验证奖励强化学习（RLVR）}] 用外部 verification function 直接产生奖励的强化学习做法。它与 preference-based RL 的差别在于：奖励 grounding 来自可检验答案或规则，而不是偏好 judge。

\item[\textbf{预算强制（budget forcing）}] 明确控制 test-time budget，使模型在更大预算下继续输出 reasoning 或继续探索的做法。它是预算分配策略，不保证自动得到更好 reasoning。

\item[\textbf{自奖励语言模型（self-rewarding language model）}] 同时承担 acting 与 judging 的模型，使系统可以从自身生成中扩展 preference data。但它仍然需要可靠 criteria 与外部 grounding，不能被理解为“随便给自己打分”。

\item[\textbf{迭代推理偏好优化（IRPO）}] 针对 reasoning task 构造 reasoning-aware preference pairs，再迭代优化策略的做法。其关键前提通常是 final answer 可验证，或至少能得到足够可靠的 negative examples。

\item[\textbf{元奖励学习（Meta-Rewarding）}] 把 judge 本身当成持续学习对象，提升评价器而不只提升生成器。它解决的是 evaluator 落后于 actor 的系统瓶颈。

\item[\textbf{评估规划器（EvalPlanner）}] 先规划评估步骤、再执行评分的 judge。它把 evaluation task 也视为 reasoning task，而不是单步打分。
\end{description}

\section{工具、代码与环境交互}

\begin{description}
\item[\textbf{工具使用（tool use）}] 智能体调用外部可执行能力，如搜索、执行、浏览、检索或验证。它是广义行为，既可以通过 function calling 实现，也可以通过更自由的 agent-computer interface 实现。

\item[\textbf{函数调用（function calling）}] 结构化 API 调用接口，用于把模型输出约束为预定义参数模式。它是 tool use 的一种接口，而不是全部 tool use。

\item[\textbf{工作流（workflow）}] 把 prompts、tools、control flow、evaluation、repair 与 guardrails 组织成可重复执行的多步流水线。书中保留该词给\textbf{系统级编排}，不用来指代单个工具调用。

\item[\textbf{编码智能体（coding agent）}] 在真实代码环境中执行定位、编辑、运行、调试、验证等动作的智能体。它强调 environment feedback，而不是单次代码补全。

\item[\textbf{智能体-计算机接口（agent-computer interface）}] 暴露给智能体的 observation/action surface，包括工具集合、反馈形状、动作粒度与 guardrails。良好的 agent-computer interface 会显著影响 agent 上限。

\item[\textbf{程序式控制（procedural control）}] 外层控制流由人工编写，模型只在指定节点被调用的系统设计。它通常更稳定、更容易调试。

\item[\textbf{动态控制（dynamic control）}] 智能体根据中间反馈在线决定下一步 action 或 tool call 的系统设计。它更灵活，但也更依赖 evaluator、sandbox 和 budget discipline。

\item[\textbf{变体分析（variant analysis）}] 从已知漏洞、补丁或 bug pattern 出发，围绕相邻代码与相似路径搜索相关问题的分析范式。它是高价值 vulnerability detection agent 的可信 framing。

\item[\textbf{前 k 命中率（pass@k）}] 在前 $k$ 个候选程序中至少有一个正确的概率。它衡量采样正确率与多样性，但不自动代表真实产品或 agent workflow 表现。
\end{description}

\section{多模态与交互式智能体}

\begin{description}
\item[\textbf{多模态智能体（multimodal agent）}] 同时消费或作用于文本、图像、界面状态、视频等多种模态的智能体。它是总类，不等于某个特定 environment。

\item[\textbf{Web 智能体（web agent）}] 主要在浏览器或网站环境中执行任务的智能体。

\item[\textbf{图形界面智能体（GUI agent）}] 在可视界面元素上完成定位与交互的智能体。它强调 visual grounding 与 action grounding。

\item[\textbf{操作系统智能体（OS agent）}] 跨桌面应用、文件系统或系统级资源执行任务的智能体。

\item[\textbf{计算机使用智能体（computer-use agent）}] 覆盖多 app、多窗口、真实 OS 环境的广义交互式智能体。该术语比 web agent 或 GUI agent 更宽。

\item[\textbf{视觉定位（visual grounding）}] 把自然语言目标对齐到截图、页面或 GUI 中具体可交互元素与布局线索的能力。它回答的是“知道该点哪里、该看哪里”。

\item[\textbf{基于执行的评测（execution-based evaluation）}] 通过环境状态、脚本检查或任务完成条件来判断成功，而不是只比较输出文本表面相似度。

\item[\textbf{引导式回放（guided replay）}] 从教程文本、演示或轨迹中恢复可执行交互序列，以生成可用于 agent training 的数据。

\item[\textbf{思维-动作链（CoTA, Chains-of-Thought-and-Action）}] 同时包含 reasoning trace 与 action trace 的训练数据或建模对象。它与只生成 thought 的 CoT 明确区分。

\item[\textbf{时序编码器（temporal encoder）}] 把长视频或长观测序列压缩成少量 latent tokens 的模块，用于把 memory bottleneck 变成可计算问题。

\item[\textbf{合成智能体任务（synthetic agentic tasks）}] 自动生成并自动验证的 agent 任务与轨迹，用于扩展训练分布。

\item[\textbf{规划-编排-学习（Plan-Seq-Learn）}] 用语言规划、场景编排与局部控制学习分层解决长程 manipulation 的方法。书中用它代表一种\textbf{层级式 agent 设计观}。
\end{description}

\section{形式化数学、证明与发现}

\begin{description}
\item[\textbf{形式化数学（formal mathematics）}] 用 proof assistant 可检查的语言表达定理、定义和证明的数学工作形态。

\item[\textbf{形式推理（formal reasoning）}] 在显式形式系统中进行、并受 machine-checkable 状态约束的推理。

\item[\textbf{形式化（formalization）}] 把非形式对象改写为形式对象的总称，可作用于 theorem、proof、program spec 或 abstraction。

\item[\textbf{形式表示（formal representation）}] 问题、概念或状态的 machine-checkable 形式表达。它是更宽的总类，不只限于 theorem statement。

\item[\textbf{形式规范（formal specification）}] 用精确定义表达 theorem、程序或任务要求的正式陈述。它回答“到底要证明或满足什么”。

\item[\textbf{自动形式化（autoformalization）}] 自动把 informal mathematics 或 natural-language statement 转换成 formal expression 的过程。

\item[\textbf{验证（verification）}] 检查候选输出是否满足形式规则、规格或 proof checker 的要求。它是\textbf{判定}动作，不等于 search。

\item[\textbf{定理证明（theorem proving）}] 在形式系统中构造一条合法证明的任务。

\item[\textbf{证明搜索（proof search）}] 对 tactics、premises、proof states 或 proof trees 进行算法探索的过程。它是 theorem proving 的内层过程，不是 theorem proving 的全部定义。

\item[\textbf{定理等价性检查（theorem equivalence checking）}] 判断两个 formal theorem 是否语义等价。

\item[\textbf{锤式自动证明系统（hammer）}] 连接 proof assistant 与外部 ATP 的系统，一般包含 premise selection、格式转换、ATP 调用与 proof reconstruction。

\item[\textbf{前提选择（premise selection）}] 从大规模 library 中选出与当前 goal 最相关的 lemmas、definitions 或 facts。

\item[\textbf{思维增强证明（thought-augmented proving）}] 在 tactic 之前显式生成 informal thought，用来为低层 proof search 提供偏置。

\item[\textbf{草图引导证明（sketch-guided proving）}] 先给出高层 proof sketch，再让 formal prover 在此约束下完成低层证明。

\item[\textbf{发现型智能体（discovery agent）}] 能提出 hypothesis、组织 search、利用 feedback 并持续抽象出 reusable concepts 的智能体。

\item[\textbf{符号回归（symbolic regression）}] 在数据上搜索紧凑、可解释的公式或程序化假设的任务。

\item[\textbf{概念库（concept library）}] 从成功 hypothesis、proof 或抽象中积累出的可复用概念集合，用于重构后续搜索空间。

\item[\textbf{视觉概念库（visual concept library）}] 针对视觉任务演化出的概念库，服务于解释、搜索与组合泛化。
\end{description}

\section{安全边界与系统防御}

\begin{description}
\item[\textbf{智能体式 AI 安全与安全防护（agentic AI safety/security）}] 面向可调用工具、可持续运行、可持久写入 memory、并能对外部环境执行动作的 agent 系统安全问题。它不等于传统的单轮 LLM safety。

\item[\textbf{直接提示注入（direct prompt injection）}] 恶意指令直接进入 prompt 空间，与 system 或 user 指令竞争。

\item[\textbf{间接提示注入（indirect prompt injection）}] 恶意内容埋在网页、邮件、文档、知识库等外部数据源中，随后被系统读取并注入上下文。

\item[\textbf{记忆投毒（memory poisoning）}] 攻击者污染长期记忆、RAG knowledge base 或其它可持续检索存储，使恶意内容跨任务反复被召回。

\item[\textbf{最小权限（least privilege）}] 组件在当前步骤只获得完成任务所需的最小能力。它是权限分配原则。

\item[\textbf{权限分离（privilege separation）}] 把高权限能力从高风险逻辑中拆分出来，减少单点失陷后的横向扩张。它是架构隔离原则，而不是权限大小本身。

\item[\textbf{信息流跟踪（information flow tracking, IFT）}] 跟踪敏感信息在组件、tool、memory 与输出之间的传播路径，用于阻止越权泄漏。
\end{description}

\section{全书易混术语的统一裁决}

\begin{itemize}
\item \textbf{Reasoning / planning / search / inference-time computation}：\emph{inference-time computation} 是预算；\emph{reasoning} 是中间推理过程；\emph{search} 是显式多候选探索；\emph{planning} 是面向未来动作的结构化决策。
\item \textbf{Post-training family}：\emph{post-training} 是总类；\emph{RLHF} 是反馈家族；\emph{DPO} 和 \emph{PPO} 是不同优化实现；\emph{RLVR} 则依赖外部 verifier，而不是 preference judge。
\item \textbf{Tool stack}：\emph{tool use} 是广义行为；\emph{function calling} 是一种结构化接口；\emph{workflow} 是把工具调用、控制流、验证和 repair 串起来的系统。
\item \textbf{Formal stack}：\emph{formalization} 是总称；\emph{formal specification} 是目标陈述；\emph{autoformalization} 是自动化形式化；\emph{verification} 负责检查；\emph{theorem proving} 负责构造证明；\emph{proof search} 是 theorem proving 的搜索内核。
\item \textbf{Security stack}：\emph{least privilege} 是权限配置原则；\emph{privilege separation} 是架构拆分原则；\emph{prompt injection} 与 \emph{memory poisoning} 是攻击方式；\emph{IFT} 是监测和约束机制。
\end{itemize}
